\section{用户需求分析}
本章从智能家居这一典型边缘部署场景出发，先说明市场规模与政策导向，再分析端侧推理的使用需求与工程约束，并补充\DrivingExtensionDirection{}方向，说明本项目的需求响应与应用延伸。本章依据用于界定问题，第五章实现方案据此展开。

\subsection{市场与政策背景}
智能家居是端侧智能计算最具规模的应用场景之一。2025 年，我国智能家居产业规模已突破万亿元，形成多元主体参与、跨界融合的产业格局；在国家补贴政策推动下，全年出货量预计达 2.79 亿台，同比增长 4.6\%，说明产业在总量扩张之外还存在结构性增长空间。\cite{smarthomescale,smarthomeshipment}

政策层面同样指向智能化升级。2026 年 7 月，国务院印发《扩大消费“十五五”规划》，提出支持家居家电消费智能化升级、发展新一代智能终端、促进智能家居设备互联互通，推动单品智能向全屋智能发展。2026 年 6 月，商务部等 8 部门发布《关于加快“人工智能+消费”发展的实施意见》，提出研究将智能家居及服务应用纳入“好房子”建设指南，加快人工智能技术在“好房子”建设中的应用。\cite{smarthomepolicy,aiconsumptionpolicy}

上述规模与政策表明，家庭设备正在从单一联网走向具备本地感知与判断能力的智能终端。设备侧是否需要具备推理能力、以何种方式部署模型，成为影响这一进程能否从“单品智能”走向“全屋智能”的工程问题。\cite{challenge}

\subsection{边缘场景的使用需求}
智能家居、安防、可穿戴与小型工业现场，都需要把现场的数据尽快转化为判断或控制动作。与传统把原始数据上传云端处理的方式相比，将模型推理放在设备附近具有三方面的直接需求：

\begin{enumerate}
\item \textbf{低时延的本地响应。}门铃、跌倒检测、危险告警等场景要求在本地即时给出结果，等待一次云端往返会明显影响可用性。推理在设备侧完成，可以把响应时间控制在与输入采集同一量级。
\item \textbf{数据不出设备的隐私诉求。}摄像头画面、语音与环境传感器读数包含大量个人信息，用户与监管都倾向于让原始数据留在设备内，只向上层输出分类、事件或摘要结果。
\item \textbf{断网可用的可靠性。}家庭网络中断或弱网环境下，照明、门锁与告警等基础功能仍应继续工作。端侧推理不依赖外部服务器，能够在离线时保持核心功能。
\end{enumerate}

这些需求共同指向同一方向：让设备在本地具备可解释、可核对的推理能力，而不是把所有计算都交给云端。

\subsection{场景对计算系统提出的约束}
边缘设备的资源条件，决定了推理系统不能照搬数据中心的软硬件结构。行业已有的端侧实现数据，进一步说明了可用计算方案必须满足的约束。\cite{k230report,ra8p1,rvvcnn,k230pytorch,cmsisnn,tinyml,ra8scenario}

\textbf{算力有限且异构。}家庭网关、摄像头等设备通常同时具备通用处理器、向量单元与轻量矩阵加速器。以 CanMV K230 为例，其玄铁 C908 双核中大核实现 RVV 1.0（向量寄存器长度 128 位，最高主频 1.6\textasciitilde GHz），小核主频 800\textasciitilde MHz 且不实现 RVV；相比标量实现，K230 上经 RVV 优化的深度卷积算子性能提升约 1.35\textasciitilde 3.8 倍。IEEE 的相关实验也表明，RVV 向量化对加法、GEMM 等计算密集型操作尤其有效，最高可实现约 23 倍加速，且性能随向量寄存器长度扩展而提升、受内存带宽影响较小。矩阵加速方面，瑞萨 RA8P1 集成 Arm Ethos-U55 NPU（256 GOPS@500MHz，可配置至 512 GOPS），相较单独使用 Cortex-M85 主核，每秒推理次数最高可提升 35 倍。这说明把不同运算分配到合适单元，是端侧性能的关键来源。\cite{k230report,k230pytorch,rvvcnn,ra8p1}

\textbf{内存、功耗与算子级优化同等重要。}端侧存储与供电有限，模型需要量化压缩，算子实现本身也直接影响效率。CMSIS-NN 在 Cortex-M7 @216\textasciitilde MHz、CIFAR-10 CNN 上的实验显示，相对基线实现，推理时间由 456.4\textasciitilde ms 降至 99.1\textasciitilde ms（吞吐提升 4.6 倍，能效提升 4.9 倍），其中卷积层提升 4.6 倍、池化层 5.4 倍、ReLU 层 2.6 倍。这提示推理系统不仅要选对硬件，还要有可复用的算子库和逐层的数值与性能对照手段。\cite{cmsisnn}

\textbf{轻量任务与丰富场景并存。}TinyML 关键词识别在 Cortex-M4F @64\textasciitilde MHz 上推理时间约 60\textasciitilde ms，模型参数量约 1000，说明小模型可以在极低资源下完成实时任务；而瑞萨 RA8 系列 MCU 支持单芯片完成驾驶员监测、语音降噪、电弧故障检测与电机故障诊断，无需外挂 DRAM，并提供视觉 AI、物体识别、异常检测等成熟 Demo。同一类端侧平台需要同时覆盖轻量任务与较复杂的视觉场景。\cite{tinyml,ra8scenario}

由此形成四类具体约束：\textbf{（1）算力有限且异构}，需要按运算结构把矩阵乘、偏置、激活与数据整理分配到合适的计算单元；\textbf{（2）内存与功耗紧张}，模型需量化压缩，运行时复用固定缓冲区，避免频繁动态分配；\textbf{（3）操作系统多样}，同一设备可能同时涉及实时系统与应用型系统，推理核心需要在两类系统下保持一致；\textbf{（4）开发与调试门槛高}，不同后端的数据布局、调用方式与数值规则存在差异，开发者难以判断模型能否部署、结果是否正确。

其中第三点在已有平台中已有对照：K230 标准开发套件中，大核运行 RT-Thread 执行机器学习推理，小核运行 Linux 作为访问板卡的便利通道。\cite{k230report}本项目在同一硬件原型上分别适配\RealTimeSystem{}与\LinuxDistribution{}，共用推理核心源码，通过重启切换镜像验证两条运行路径。\cite{technicalreport}

\subsection{本项目的需求响应}
\ProjectName{}以统一算子语义、可核验的多后端执行和静态执行计划回应上述需求。系统在\ProcessorCore{}、\VectorUnit{}与\AcceleratorName{}组成的异构结构上，由电脑端完成模型检查与后端分配，板端依次执行并输出带样本标识的结果；同一推理核心源码在\RealTimeSystem{}与\LinuxDistribution{}下复用，通过重启切换镜像完成适配。相关设计详见第四章。

上述行业数据同时说明，端侧性能来自向量化、算子库与加速器分工的共同作用，而非单一硬件指标。因此本项目把验证重点放在逐层数值对照与后端比较上：以独立参考实现确认整数结果一致，再分别记录各后端的适用范围。首期以受限小模型和静态输入验证推理链路的连通性，为后续接入摄像头、传感器与通信接口保留明确的输入输出边界。端侧推理的时延、隐私与可靠性需求能否满足，将以逐层数值对照、两种系统下的重复运行以及真实输入联调的结果作为依据；具体性能与能耗指标在工程验证后另行记录。\cite{challenge,technicalreport,rvvcnn,cmsisnn}

\par\textbf{\DrivingExtensionDirection{}需求。}车内视觉感知可服务于\DrivingPilotTask{}，例如从图像中识别\DrivingPilotStates{}等状态；已有车载边缘 AI 方案采用本地处理支持此类任务。\cite{cabinai}本项目可进一步研究连续图像输入下的推理时延、状态变化记录和原始图像本地处理，通过统一张量接口与算子部署流程衔接车载视觉模型。
